\section{Model Optimization and Deployment}
\label{sec:deployment}

Deploying the trained GestureTCN on edge devices requires reducing both
model size and inference latency. We apply a three-step optimization
pipeline: structured pruning, ONNX export, and INT8 quantization.

\subsection{Structured Pruning}

Unlike unstructured pruning (which zeroes individual weights but
maintains tensor dimensions), \emph{structured pruning} removes entire
convolutional channels, producing a genuinely smaller
model~\cite{li2017pruning}. We apply a uniform channel reduction ratio
of 30\%, with channel counts rounded to multiples of~8 for SIMD
alignment:

\begin{table}[ht]
\centering
\caption{Channel dimensions before and after structured pruning.}
\label{tab:pruning}
\begin{tabular}{@{}lcc@{}}
\toprule
\textbf{Component} & \textbf{Original} & \textbf{Pruned} \\
\midrule
Stem channels & 48 & 32 \\
Mid-block channels & 48 & 32 \\
Output channels & 64 & 48 \\
Head hidden units & 32 & 24 \\
\midrule
Total parameters & 87,077 & 45,877 \\
Model size (FP32) & 0.34\,MB & 0.18\,MB \\
Compression ratio & 1.0$\times$ & 1.9$\times$ \\
\bottomrule
\end{tabular}
\end{table}

The pruned model is a new instance of \texttt{GestureTCN} with reduced
channel dimensions. It is fine-tuned for 100~epochs using the
same dataset and a lower learning rate ($10^{-3}$), achieving a test
accuracy of 89.66\%.

\subsection{ONNX Export}

We export the pruned model to ONNX (Open Neural Network
Exchange)~\cite{onnxruntime} format using PyTorch's tracing-based
exporter. The export specification includes a dynamic batch axis to accommodate variable batch sizes at inference time, opset version 18 to ensure compatibility with modern ONNX Runtime versions, and named input and output tensors where the input tensor is designated as \texttt{``input''} with shape $[B, 144, 30]$ and the output tensor is designated as \texttt{``logits''} with shape $[B, 5]$.
Model validity is verified using \texttt{onnx.checker.check\_model()}.

\subsection{INT8 Static Quantization}

We apply static quantization via ONNX
Runtime~\cite{jacob2018quantization}, converting FP32 weights to
signed 8-bit integers through a three-phase process. In the calibration phase, real training data is fed through the FP32 ONNX model to observe the range of activation values at each layer. Subsequently, in the scale computation phase, per-tensor scale and zero-point values are computed to map FP32 ranges to INT8 according to the formula $x_{\text{int8}} = \text{round}(x_{\text{fp32}} / s) + z$. Finally, in the QDQ insertion phase, Quantize-Dequantize (QDQ) nodes are inserted into the computation graph, enabling the runtime to fuse quantized operations for optimal performance.

\paragraph{Quantization Benefits.}
The primary benefit of INT8 quantization for this model is \emph{inference acceleration} rather than size reduction. On mobile CPUs and NPUs, INT8 operations are typically 2--4$\times$ faster than FP32 operations due to specialized hardware support and reduced memory bandwidth. While quantization theoretically reduces weight storage by 4$\times$ (32-bit to 8-bit), the QDQ format adds metadata overhead for each quantized operation. For small models ($<$100\,KB), this overhead can offset weight compression gains---our quantized model is 0.14\,MB compared to the FP32 ONNX model's 0.09\,MB. This phenomenon is well-documented for compact neural networks where the fixed overhead of quantization metadata dominates the variable savings from weight compression.

\subsection{Deployment Configuration}

A \texttt{config.json} file is generated alongside the ONNX model,
containing all information required for inference on edge devices:
class names, sequence length, feature dimension, normalization
statistics ($\boldsymbol\mu$ and $\boldsymbol\sigma$ vectors of
dimension~144), landmark topology (finger chains, joint pairs,
tip/base indices), and pruned channel dimensions. This ensures
complete self-containment of the deployed artifact.

\subsection{Optimization Summary}

Table~\ref{tab:optimization_summary} summarizes the full optimization
pipeline.

\begin{table}[ht]
\centering
\caption{Model optimization pipeline summary.}
\label{tab:optimization_summary}
\begin{tabular}{@{}lcccc@{}}
\toprule
\textbf{Stage} & \textbf{Params} & \textbf{Size} & \textbf{Accuracy} & \textbf{Format} \\
\midrule
Original (FP32) & 87,077 & 0.34\,MB & 89.66\% & PyTorch \\
Pruned (FP32) & 45,877 & 0.18\,MB & 89.66\% & PyTorch \\
Original ONNX (FP32) & 87,077 & 0.09\,MB & 89.66\% & ONNX \\
Pruned ONNX (FP32) & 45,877 & 0.09\,MB & 89.66\% & ONNX \\
Pruned + Quantized (INT8) & 45,877 & 0.14\,MB & 89.66\% & ONNX \\
\bottomrule
\end{tabular}
\end{table}

The ONNX format produces more compact files than PyTorch's serialization format. Notably, the INT8 quantized model is slightly larger than the FP32 ONNX models due to QDQ node overhead, as discussed above. For deployment, we select the quantized model primarily for its faster inference speed on mobile hardware rather than size reduction.
